% Options for packages loaded elsewhere
\PassOptionsToPackage{unicode}{hyperref}
\PassOptionsToPackage{hyphens}{url}
\documentclass[
  american,
  11pt,
  letterpaper,
]{article}
\usepackage{xcolor}
\usepackage[margin=1in]{geometry}
\usepackage{amsmath,amssymb}
\setcounter{secnumdepth}{-\maxdimen} % remove section numbering
\usepackage{iftex}
\ifPDFTeX
  \usepackage[T1]{fontenc}
  \usepackage[utf8]{inputenc}
  \usepackage{textcomp} % provide euro and other symbols
\else % if luatex or xetex
  \usepackage{unicode-math} % this also loads fontspec
  \defaultfontfeatures{Scale=MatchLowercase}
  \defaultfontfeatures[\rmfamily]{Ligatures=TeX,Scale=1}
\fi
\usepackage{lmodern}
\ifPDFTeX\else
  % xetex/luatex font selection
\fi
% Use upquote if available, for straight quotes in verbatim environments
\IfFileExists{upquote.sty}{\usepackage{upquote}}{}
\IfFileExists{microtype.sty}{% use microtype if available
  \usepackage[]{microtype}
  \UseMicrotypeSet[protrusion]{basicmath} % disable protrusion for tt fonts
}{}
\usepackage{setspace}
\makeatletter
\@ifundefined{KOMAClassName}{% if non-KOMA class
  \IfFileExists{parskip.sty}{%
    \usepackage{parskip}
  }{% else
    \setlength{\parindent}{0pt}
    \setlength{\parskip}{6pt plus 2pt minus 1pt}}
}{% if KOMA class
  \KOMAoptions{parskip=half}}
\makeatother
\usepackage{graphicx}
\makeatletter
\newsavebox\pandoc@box
\newcommand*\pandocbounded[1]{% scales image to fit in text height/width
  \sbox\pandoc@box{#1}%
  \Gscale@div\@tempa{\textheight}{\dimexpr\ht\pandoc@box+\dp\pandoc@box\relax}%
  \Gscale@div\@tempb{\linewidth}{\wd\pandoc@box}%
  \ifdim\@tempb\p@<\@tempa\p@\let\@tempa\@tempb\fi% select the smaller of both
  \ifdim\@tempa\p@<\p@\scalebox{\@tempa}{\usebox\pandoc@box}%
  \else\usebox{\pandoc@box}%
  \fi%
}
% Set default figure placement to htbp
\def\fps@figure{htbp}
\makeatother
\ifLuaTeX
\usepackage[bidi=basic,shorthands=off,]{babel}
\else
\usepackage[bidi=default,shorthands=off,]{babel}
\fi
\ifLuaTeX
  \usepackage{selnolig} % disable illegal ligatures
\fi
\setlength{\emergencystretch}{3em} % prevent overfull lines
\providecommand{\tightlist}{%
  \setlength{\itemsep}{0pt}\setlength{\parskip}{0pt}}
\usepackage{booktabs}
\usepackage{longtable}
\usepackage{array}
\usepackage{float}
\usepackage{caption}
\usepackage{fancyhdr}
\pagestyle{fancy}
\fancyhf{}
\fancyhead[L]{Piracy and PIPA/SOPA}
\fancyhead[R]{Decision Making Under Uncertainty}
\fancyfoot[C]{\thepage}
\setlength{\headheight}{14pt}
\usepackage{etoolbox}
\makeatletter
\renewcommand{\maketitle}{
  \begin{titlepage}
  \thispagestyle{empty}
  \vspace*{\fill}

  \begin{center}
    {\LARGE \@title \par}
    \vspace{1em}

    {\large
    \lineskip .5em
    \begin{tabular}[t]{c}
    \@author
    \end{tabular}\par}

    \vspace{1em}
    {\large \@date \par}

    \vspace{2em}
    \begin{minipage}{0.7\textwidth}
    \centering
    \textbf{Abstract}\\[0.5em]
    We need to put something here which gives an abstract description of what we did and why it matters
    \end{minipage}
  \end{center}

  \vspace*{\fill}
  \end{titlepage}
}

\makeatother
\pretocmd{\tableofcontents}{\thispagestyle{empty}}{}{}
\apptocmd{\tableofcontents}{\newpage}{}{}
\usepackage{bookmark}
\IfFileExists{xurl.sty}{\usepackage{xurl}}{} % add URL line breaks if available
\urlstyle{same}
\hypersetup{
  pdftitle={Piracy and PIPA/SOPA},
  pdfauthor={Christian Auman; Nate Miller; Nathan; Simran Gautam},
  pdflang={en-US},
  hidelinks,
  pdfcreator={LaTeX via pandoc}}

\title{Piracy and PIPA/SOPA}
\usepackage{etoolbox}
\makeatletter
\providecommand{\subtitle}[1]{% add subtitle to \maketitle
  \apptocmd{\@title}{\par {\large #1 \par}}{}{}
}
\makeatother
\subtitle{This is our final project for Decision Making Under
Uncertainty}
\author{Christian Auman \and Nate Miller \and Nathan \and Simran Gautam}
\date{April 25, 2026}

\begin{document}
\maketitle

{
\setcounter{tocdepth}{2}
\tableofcontents
}
\setstretch{1.5}
\section{Introduction}\label{introduction}

Give a brief introductory paragraph on the project.

\subsection{Project Motivation}\label{project-motivation}

\textbf{Questions to Answer:}

\begin{itemize}
\tightlist
\item
  Why do we care?
\item
  What are we doing?
\item
  Briefly, how are we solving it?
\item
  What is the data?
\item
  Where does it come from?
\end{itemize}

\subsection{Data Source and Variables}\label{data-source-and-variables}

\begin{quote}
More in-depth overview about the data
\end{quote}

\begin{itemize}
\tightlist
\item
  Population vs Sample
\item
  Size:

  \begin{itemize}
  \tightlist
  \item
    How many observations?
  \item
    How many variables?
  \end{itemize}
\item
  Variables - describe each variable in the dataset (more than just it's
  type)
\end{itemize}

\section{Research Questions and
Hypothesis}\label{research-questions-and-hypothesis}

Give a brief introductory paragraph on this section.

\subsection{Research Question}\label{research-question}

\begin{itemize}
\tightlist
\item
  Define our primary question that our analysis will answer
\item
  Describe how we came up with this question
\end{itemize}

\subsection{Our Hypothesis}\label{our-hypothesis}

\begin{itemize}
\tightlist
\item
  State null \& alternative hypotheses in english
\item
  Then describe them in statistics language (i.e.~symbols)
\item
  List 2-3 supporting questions we will investigate
\end{itemize}

\section{Main Analyses}\label{main-analyses}

Briefly describe this section.

\subsection{Exploratory Data Analysis}\label{exploratory-data-analysis}

Give an overview of the plots we use \& statistics we used and relate
how they answer our research questions.

Describe the plot and how it helps answer a research question

\textbf{PLOT 1} - PLOT 1 figcaption

Describe the plot and how it helps answer a research question

\textbf{PLOT 2} - PLOT 2 figcaption

Describe the plot and how it helps answer a research question

\textbf{PLOT 3} - PLOT 3 figcaption

\subsection{Statistical Methods}\label{statistical-methods}

List the statistical analysis methods we used to answer our research
question.

\begin{itemize}
\tightlist
\item
  Go through each method in separate sections

  \begin{itemize}
  \tightlist
  \item
    describe how each method answers one of our research question
  \item
    explain the results of the particular analysis
  \item
    define the analysis in proper statistics notation (i.e.~with
    symbols)
  \end{itemize}
\end{itemize}

\section{Discussion (Interpretation, uncertainty,
assumptions)}\label{discussion-interpretation-uncertainty-assumptions}

\emph{can merge this section into conclusion if we need to shorten the
report}

Briefly interpret the results beyond the raw statistics.

\begin{itemize}
\tightlist
\item
  What do the results imply?
\item
  Were the hypotheses supported?
\item
  How should we understand the uncertainty?
\item
  What assumptions did the analysis rely on?
\end{itemize}

\section{Conclusion (final answer, limitations, future work, why it
matters)}\label{conclusion-final-answer-limitations-future-work-why-it-matters}

Briefly summarize the project and directly answer the main research
question.

\subsection{Summary of Findings}\label{summary-of-findings}

\begin{itemize}
\tightlist
\item
  Restate our finding: What did the tests conclude (in plain english)
\item
  Answer our research question directly
\item
  Supporting findings: What did the exploratory analyses reveal
\end{itemize}

\subsection{Limitations and Future
Work}\label{limitations-and-future-work}

\begin{itemize}
\tightlist
\item
  Limitations: What are the weaknesses of our analysis? (most probably
  the correlation coefficient is not strong)
\item
  Unanswered Questions: What remains unknown?
\item
  Future research section
\item
  Real-world meaning: So what? Why should anyone care about our
  findings?
\end{itemize}

\end{document}
